\begin{abstract}

This paper applies modern causal inference methods to examine how institutional structures, policy environments, and contextual conditions shape behavioral responses in policing, corrections, and consumer markets. The first chapter investigates how law enforcement behavior shifts in the aftermath of police-shooting fatalities by combining tens of millions of traffic stop records with fatal encounters timing data. The second chapter analyzes whether reducing communication barriers improves safety outcomes for incarcerated individuals by leveraging the 2014 Federal Communications Commission (FCC) mandate to cap prison phone rates. The third chapter provides the first national assessment of heterogeneity in gasoline price elasticity across U.S. states by examining newly developed state-level elasticity estimates in conjunction with data on state environments, road infrastructure, and population characteristics. \\

In Chapter 1, I analyze traffic stop and search patterns following lethal events, providing causal estimates that policing outcome volumes drop significantly following a death event, a finding consistent with decreases in policing efforts. Exploring event heterogeneity across public scrutiny levels, I find further suggestive evidence that these effort decreases follow from low scrutiny, but under high scrutiny it is search success rates that decrease. \\

In Chapter 2, I combine manually-collect data from prison telecommunications contracts with administrative records on inmate sexual assault reports, obtaining a null result for the causal effect of the FCC rate cap on prisoner well being, suggesting that while easing financial constraints on communication may enhance inmate welfare along other dimensions, phone rate reforms alone do not appear to meaningfully influence safety or reporting outcomes. \\

In Chapter 3, a joint paper with Michael Bates \& Seolah Kim, we provide the first empirical assessment of heterogeneity in gas price elasticity across the United States by combining new elasticity estimates with data on state geographies, climate, road infrastructure, and population measures. Our findings show that both larger states and those with more favorable weather have relatively inelastic demand, while states with more dense populations, larger border populations, and more urbanized road networks exhibit higher elasticity, carrying important implications for future policy development. 

\end{abstract}